% GENERATED FILE. Edit canonical lessons, then run npm run generate:books.
% canonical-source-hash: fnv1a64:076dc4e828ba0275
% canonical-lessons: TA-C167-totar, TA-C167-iyakku, TA-C167-ottu, TA-C167-tonku, TA-C167-uruttu

\chapter{To continue, To operate, To stick, To hang, To roll}
\label{ch:ta-to-continue-to-operate-to-stick-to-hang-to-roll}
\hlchaptermodality{167}

\begin{chapteropening}
\textbf{By the end of this chapter:} \emph{I can say to continue, to operate, to stick, to hang and to roll.}

Complete the last lesson of chapter 167: I can say to continue, to operate, to stick, to hang and to roll.
\end{chapteropening}

\section[\texorpdfstring{\ta{தொடர்} (toṭar)}{தொடர் (toṭar)}]{\ta{தொடர்} (toṭar) — to continue}

\label{lesson:TA-C167-totar}



\begin{quote}
Before the new one: say the Tamil for to make peace, then the Tamil for to consider.
\end{quote}

\subsection*{You'll want to know: \ta{தொடர்}}
\textbf{\ta{தொடர்}} — \emph{toṭar} — \textquotedblleft{}to continue\textquotedblright{}.

\begin{grammarlens}[title={What you've built}]
One more word for this chapter.
\end{grammarlens}

\subsection*{Your turn}
\begin{itemize}
\raggedright
  \item \emph{Say it:} \emph{toṭar}
  \item \emph{Say it:} \emph{toṭar}, once more
  \item \emph{Say it:} say \emph{toṭar}
  \item \emph{Recall:} say the Tamil for to make peace, then the Tamil for to consider, then say \emph{toṭar} again
\end{itemize}

\begin{tcolorbox}[breakable,colback=teal!4,colframe=teal!35!black,title={Before you move on}]
What does \ta{தொடர்} mean? (\textquotedblleft{}To continue\textquotedblright{}.) Say it once more.
\end{tcolorbox}
\section[\texorpdfstring{\ta{இயக்கு} (iyakku)}{இயக்கு (iyakku)}]{\ta{இயக்கு} (iyakku) — to operate, to run (a machine)}

\label{lesson:TA-C167-iyakku}



\begin{quote}
Before the new one: say the Tamil for to consider, then the Tamil for to continue.
\end{quote}

\subsection*{You'll want to know: \ta{இயக்கு}}
\textbf{\ta{இயக்கு}} — \emph{iyakku} — \textquotedblleft{}to operate, to run (a machine)\textquotedblright{}.

\begin{grammarlens}[title={What you've built}]
One more word for this chapter.
\end{grammarlens}

\subsection*{Your turn}
\begin{itemize}
\raggedright
  \item \emph{Say it:} \emph{iyakku}
  \item \emph{Say it:} \emph{iyakku}, once more
  \item \emph{Say it:} say \emph{iyakku}
  \item \emph{Recall:} say the Tamil for to consider, then the Tamil for to continue, then say \emph{iyakku} again
\end{itemize}

\begin{tcolorbox}[breakable,colback=teal!4,colframe=teal!35!black,title={Before you move on}]
What does \ta{இயக்கு} mean? (\textquotedblleft{}To operate, to run (a machine)\textquotedblright{}.) Say it once more.
\end{tcolorbox}
\section[\texorpdfstring{\ta{ஒட்டு} (oṭṭu)}{ஒட்டு (oṭṭu)}]{\ta{ஒட்டு} (oṭṭu) — to stick, to paste}

\label{lesson:TA-C167-ottu}



\begin{quote}
Before the new one: say the Tamil for to continue, then the Tamil for to operate.
\end{quote}

\subsection*{You'll want to know: \ta{ஒட்டு}}
\textbf{\ta{ஒட்டு}} — \emph{oṭṭu} — \textquotedblleft{}to stick, to paste\textquotedblright{}.

\begin{grammarlens}[title={What you've built}]
One more word for this chapter.
\end{grammarlens}

\subsection*{Your turn}
\begin{itemize}
\raggedright
  \item \emph{Say it:} \emph{oṭṭu}
  \item \emph{Say it:} \emph{oṭṭu}, once more
  \item \emph{Say it:} say \emph{oṭṭu}
  \item \emph{Recall:} say the Tamil for to continue, then the Tamil for to operate, then say \emph{oṭṭu} again
\end{itemize}

\begin{tcolorbox}[breakable,colback=teal!4,colframe=teal!35!black,title={Before you move on}]
What does \ta{ஒட்டு} mean? (\textquotedblleft{}To stick, to paste\textquotedblright{}.) Say it once more.
\end{tcolorbox}
\section[\texorpdfstring{\ta{தொங்கு} (toṅku)}{தொங்கு (toṅku)}]{\ta{தொங்கு} (toṅku) — to hang (intransitive), to dangle}

\label{lesson:TA-C167-tonku}



\begin{quote}
Before the new one: say the Tamil for to operate, then the Tamil for to stick.
\end{quote}

\subsection*{You'll want to know: \ta{தொங்கு}}
\textbf{\ta{தொங்கு}} — \emph{toṅku} — \textquotedblleft{}to hang (intransitive), to dangle\textquotedblright{}.

\begin{grammarlens}[title={What you've built}]
One more word for this chapter.
\end{grammarlens}

\subsection*{Your turn}
\begin{itemize}
\raggedright
  \item \emph{Say it:} \emph{toṅku}
  \item \emph{Say it:} \emph{toṅku}, once more
  \item \emph{Say it:} say \emph{toṅku}
  \item \emph{Recall:} say the Tamil for to operate, then the Tamil for to stick, then say \emph{toṅku} again
\end{itemize}

\begin{tcolorbox}[breakable,colback=teal!4,colframe=teal!35!black,title={Before you move on}]
What does \ta{தொங்கு} mean? (\textquotedblleft{}To hang (intransitive), to dangle\textquotedblright{}.) Say it once more.
\end{tcolorbox}
\section[\texorpdfstring{\ta{உருட்டு} (uruṭṭu)}{உருட்டு (uruṭṭu)}]{\ta{உருட்டு} (uruṭṭu) — to roll (something)}

\label{lesson:TA-C167-uruttu}



\begin{quote}
Before the new one: say the Tamil for to stick, then the Tamil for to hang.
\end{quote}

\subsection*{You'll want to know: \ta{உருட்டு}}
\textbf{\ta{உருட்டு}} — \emph{uruṭṭu} — \textquotedblleft{}to roll (something)\textquotedblright{}.

\begin{grammarlens}[title={What you've built}]
One more word for this chapter.
\end{grammarlens}

\subsection*{Your turn}
\begin{itemize}
\raggedright
  \item \emph{Say it:} \emph{uruṭṭu}
  \item \emph{Say it:} \emph{uruṭṭu}, once more
  \item \emph{Say it:} say \emph{uruṭṭu}
  \item \emph{Recall:} say the Tamil for to stick, then the Tamil for to hang, then say \emph{uruṭṭu} again
\end{itemize}

\begin{tcolorbox}[breakable,colback=teal!4,colframe=teal!35!black,title={Before you move on}]
What does \ta{உருட்டு} mean? (\textquotedblleft{}To roll (something)\textquotedblright{}.) Say it once more.
\end{tcolorbox}
